\section{Integer certificates and quantitative arithmetic rank}
\label{sec:fd-inertia-density}

The finite trace test admits an exact integer version involving the
Gram and primitive matrices. We first realize the distinct-zero density
in these matrices, and then approximate the individual spectral
coordinates of a large isolated simple family.

\subsection{The exact integer test}
For \(T>0\), let \(n(T)\) count positive zero ordinates at most
\(T\), with multiplicity, and let \(n_{\mathrm{dist}}(T)\) count
these ordinates without multiplicity. Let
\[
 \iota_N(T)=\operatorname{ind}_+(T H_N-G_N),
\]
where \(\operatorname{ind}_+\) is the number of strictly positive
eigenvalues of a finite Hermitian matrix, counted with multiplicity.
The matrices \(G_N,H_N\) are the arithmetic matrices of
Section~\ref{sec:fd-arithmetic}.

\begin{proposition}[A finite inertia certificate]
\label{prop:fd-inertia-complete}
At every height \(T>0\) with \(\Xi(T)\ne0\), the integers
\(\iota_N(T)\) increase with \(N\), are bounded above by
\(n_{\mathrm{dist}}(T)\), and eventually equal that number.
Consequently the packet through \(T\) is simple exactly when
\[
 \iota_N(T)\ge n(T)
\]
holds for one finite \(N\). Simplicity of every nontrivial zero is
equivalent to these finite inequalities at an unbounded sequence of
heights with \(\Xi(T)\ne0\).
\end{proposition}
\begin{proof}
Congruence by \(G_N^{-1/2}\) turns \(T H_N-G_N\) into
\(T C_N-I\) on \(V_N\). Sylvester's law of inertia therefore
identifies \(\iota_N(T)\) with the number of positive eigenvalues
of \(C_N\) greater than \(1/T\). The variational principle bounds
these by the positive eigenvalues of \(C\), and makes them
nondecreasing under the nested projections. The latter eigenvalues
are \(1/\gamma\), one per distinct positive ordinate.

Only finitely many exceed \(1/T\). None equals \(1/T\), by the
boundary condition. Operator-norm convergence \(C_N\to C\)
therefore places every eigenvalue above this threshold on the correct
side for all sufficiently large \(N\); the upper variational bound
excludes additional ones. This proves eventual equality.
The counted and distinct counts agree exactly for a simple packet,
which proves the finite equivalence. An unbounded sequence of simple
packets covers every positive ordinate. Conversely, when every zero
is simple, choose unbounded zero-free contour heights and apply the
finite equivalence at each one.
\end{proof}

This test uses no high-order Taylor coefficients. In an even--odd
basis of \(V_N\), both parity dimensions equal
\(r_N=N2^N-1\). Write
\[
 G_N=\begin{pmatrix}G_+&0\\0&G_-\end{pmatrix},\qquad
 H_N=\begin{pmatrix}0&iK\\-iK^{\mathsf T}&0\end{pmatrix},
\]
with real \(K\); its sign depends only on the order of the parity
basis. Eliminating the negative definite upper block in
\(T H_N-G_N\) gives the real symmetric test
\[
 \iota_N(T)=\operatorname{ind}_+
       \bigl(T^2K^{\mathsf T}G_+^{-1}K-G_-\bigr).
\]
This is an exact Schur-complement identity, not an approximation to
its eigenvalues.

\subsection{A quantitative rank bound from separated coordinates}
\label{subsec:strong-inertia-density}

The positive inertia counts distinct spectral coordinates. A selected
coordinate may have multiplicity greater than one: its sampling weight
is positive and increases the lower quadratic-form estimate. We obtain
the following quantitative count. Put
\[
 d_{\rm MT}=\frac{1+c_{\rm MT}}2
 =\frac54-\frac1{2\sqrt2}\cot\frac1{\sqrt2}
 =0.8362503\ldots.
\]
The constant \(c_{\rm MT}\) and its admissible test approximation
were proved in Section~\ref{sec:survival}.

\begin{theorem}[Quantitative arithmetic realization of distinct zeros]
\label{thm:strong-inertia-density}
For every \(0<\eta<d_{\rm MT}\), there are constants
\(A_\eta,T_\eta>0\) such that
\[
 \operatorname{ind}_+(T H_N-G_N)\ge(d_{\rm MT}-\eta)n(T),
 \qquad N=\lceil A_\eta\log T\rceil,\quad T\ge T_\eta.
\]
The matrices use only the arithmetic Weil pairings. Their dimension is
\(O_\eta(T^{A_\eta\log2}\log T)\), and their prime terms have
\(n<e^{2N}\le e^2T^{2A_\eta}\).
More precisely, for every admissible nonnegative pair test \(g\) used
in Section~\ref{sec:survival}, the same conclusion holds with any
constant smaller than \(1-I_g/2\), where
\[
 I_g=\int_{\mathbb R}g(u)
      \left(1-\left(\frac{\sin\pi u}{\pi u}\right)^2\right)du.
\]
\end{theorem}

\begin{proof}
We first construct enough mutually separated coordinates. Fix an even
nonnegative Schwartz function \(g\) with \(g(0)=1\), Fourier support
strictly inside \((-1,1)\), and \(I_g\) sufficiently close to \(1-c_{\rm MT}\).
The approximations in the proof of
Theorem~\ref{thm:two-thirds-survival} supply such tests. Choose
\(\epsilon,\delta>0\) so that
\(g(u)\ge1-\epsilon\) for \(|u|\le\delta\). The restricted
pair-correlation formula gives
\[
 \sum_{\substack{0<\gamma_j,\gamma_k\le T\\j\ne k}}
 g\bigl(F(\gamma_j)-F(\gamma_k)\bigr)
       =(I_g+o(1))n(T),\qquad F(v)=\frac{v\log v}{2\pi},
\]
with all labels and multiplicities retained. Choose fixed
\(0<\alpha<\beta<1\), with \(\alpha\) small and \(\beta\)
close to one. For large \(T\), \(F'(v)\le\log T\) throughout
\([\alpha T,\beta T]\). Thus every unordered pair of distinct
labels in this interval at distance less than
\(h=\delta/\log T\) contributes at least \(2(1-\epsilon)\)
to the displayed ordered sum. The number of such pairs is at most
\[
 \left(\frac{I_g}{2(1-\epsilon)}+o(1)\right)n(T).
\]
In the finite graph joining precisely these close pairs, delete one
endpoint of every edge. The retained labels are mutually
\(h\)-separated. The Riemann--von Mangoldt count and suitable fixed
choices of the parameters give a set
\[
 Z_T=\{\gamma_1,\ldots,\gamma_M\}\subset[\alpha T,\beta T],
 \qquad M\ge(d_{\rm MT}-\eta/2)n(T).
\]
Its ordinates are distinct. Other occurrences, including repeated
copies at a selected ordinate, remain in the sampling measure.
Keeping \(I_g\) fixed in this argument gives every count below
\((1-I_g/2)n(T)\), which proves the final quantitative dependence
once the approximation argument below is established.

Write
\[
 Q_T(f)=\sum_{\gamma\in\Gamma}m_\gamma
             (T/\gamma-1)|\widehat f(\gamma)|^2,
 \qquad
 \|f\|_T^2=\sum_{\gamma\in\Gamma}m_\gamma
             (1+T/|\gamma|)|\widehat f(\gamma)|^2.
\]
The second expression bounds the absolute sesquilinear form of the
first. If \(\gamma_0>0\) is the least absolute ordinate, then
\(\|f\|_T^2\le(1+T/\gamma_0)q(f)\).
The unit-interval counted zero bound gives
\begin{equation}\label{eq:strong-inertia-count-sums}
 \begin{split}
 n(T)&=O(T\log(3T)),\\
 \sum_{0<|\gamma|\le T}\frac{m_\gamma}{|\gamma|}
       &=O(\log^2(3T)),\\
 \sum_{|\gamma|>T}\frac{m_\gamma}{\gamma^2}
       &=O\bigl(\log(3T)/T\bigr).
 \end{split}
\end{equation}
These follow by summing over unit intervals, with the finite initial
range absorbed in the constants.

Take \(N=\lceil A\log T\rceil\), \(a=\lfloor N/2\rfloor\), and
\[
 g_a(x)=\frac1a(1-|x|/a)_+,
 \qquad
 K_a(t)=\widehat g_a(t)
       =\left(\frac{\sin(at/2)}{at/2}\right)^2.
\]
The integral of \(g_a\) is one. Define
\[
 d_j=K_a(\gamma_j),\qquad
 f_j(x)=e^{i\gamma_jx}g_a(x)-d_jg_a(x),\qquad
 Uc=\sum_{j=1}^M c_jf_j.
\]
Every \(f_j\) has integral zero and support in \([-a,a]\).
The selected sampling matrix is exactly \(I+R-dd^*\), with
\(R_{kj}=K_a(\gamma_k-\gamma_j)\) for \(k\ne j\) and zero diagonal.
The elementary bound
\(|K_a(t)|\le\min(1,4/(a^2t^2))\) and separation give
\[
 \|R\|\le\frac{4}{a^2}\,2h^{-2}\sum_{r\ge1}r^{-2},
 \qquad
 \|d\|^2\le\frac{16M}{a^4\alpha^4T^4}=o(1).
\]
Choose \(A\) sufficiently large that the selected matrix is bounded
below by \(1/2\) for every sufficiently large \(T\). Each selected
coefficient satisfies \(T/\gamma_j-1\ge1/\beta-1\), and its
multiplicity is at least one. Its total contribution to \(Q_T(Uc)\)
is therefore at least
\begin{equation}\label{eq:strong-inertia-positive}
 \frac{1/\beta-1}{4}\|c\|_2^2.
\end{equation}
All remaining coordinates in \((0,T)\) have nonnegative
coefficients and can be omitted from a lower bound.

The negative part is confined to \(\gamma<0\) or \(\gamma>T\).
For these ordinates and all selected \(\gamma_j\),
\[
 |\gamma-\gamma_j|\ge c_0(T+|\gamma|),
 \qquad c_0=\min\{\alpha,(1-\beta)/2\}>0.
\]
For the uncorrected columns \(K_a(\gamma-\gamma_j)\), the row sum
is at most \(CM/[a^2(T+|\gamma|)^2]\). With output weight
\(m_\gamma(1+T/|\gamma|)\), the column sum is at most
\[
 \frac C{a^2}\sum_\gamma
 \frac{m_\gamma(1+T/|\gamma|)}{(T+|\gamma|)^2}
       \le\frac{C\log^2(3T)}{a^2T},
\]
by \eqref{eq:strong-inertia-count-sums}. Weighted Schur therefore
bounds the squared operator norm of the entire negative sampling
part by
\begin{equation}\label{eq:strong-inertia-negative}
 \frac{CM\log^2(3T)}{a^4T^3}
       =O\left(\frac{\log^3(3T)}{a^4T^2}\right)=o(1).
\end{equation}
Here the weighted Schur inequality follows directly from
\(|\sum_j k_{\gamma j}c_j|^2
 \le(\sum_j|k_{\gamma j}|)
        (\sum_j|k_{\gamma j}||c_j|^2)\), followed by summation
with the output weights.
The mean correction is negligible in the full norm, since
\[
 \|g_a\|_T^2\le\frac{16}{a^4}
   \sum_\gamma\frac{m_\gamma(1+T/|\gamma|)}{\gamma^4}
      \le\frac{C(1+T)}{a^4},
 \qquad \|d\|\|g_a\|_T=o(1).
\]
Equations~\eqref{eq:strong-inertia-positive} and
\eqref{eq:strong-inertia-negative} imply
\begin{equation}\label{eq:strong-inertia-margin}
 Q_T(Uc)\ge c_\beta\|c\|_2^2
\end{equation}
with fixed \(c_\beta>0\), for all large \(T\).

We also need a bound on the full coefficient map. For any output
ordinate, at most two selected points lie within distance \(h\);
the remaining distances are bounded below by successive multiples of
\(h\). Thus the unweighted row sum of
\(|K_a(\gamma-\gamma_j)|\) is \(O(1)\), including arbitrarily
close unselected ordinates. For a fixed column, the counted zeros
within distance one contribute \(O(\log T)\). Outside that interval
and with \(|\gamma|\ge\alpha T/2\), the weight
\(1+T/|\gamma|\) is bounded, and unit-distance shells contribute
\[
 O\left(\frac1{a^2}\sum_{r\ge1}
                    \frac{\log(T+r)}{r^2}\right)=O(\log T).
\]
The remaining range \(|\gamma|<\alpha T/2\) contributes
\(O(\log^2(3T)/(a^2T))\), again by
\eqref{eq:strong-inertia-count-sums}. Weighted Schur and the mean
correction consequently give
\begin{equation}\label{eq:strong-inertia-fullnorm}
 \|U\|_{\mathbb C^M\to\|\cdot\|_T}=O(\sqrt{\log T}).
\end{equation}

Finally interpolate into the specified arithmetic space. Let
\(b=2^{-N}\), let \(I_b\) be nodal PL interpolation, and set
\[
 v_j=I_bf_j-g_a\int I_bf_j,\qquad U_Nc=\sum_jc_jv_j.
\]
The knots \(-a,0,a\) belong to the mesh, and \(g_a\) is an exact
mesh function. Hence \(v_j\in V_N\), with integral exactly zero.
Within the mesh cells,
\(\|f_j''\|_2\le CT^2/\sqrt a\). The derivative of a cell
interpolant is the average derivative; the one-dimensional Poincar\'e
inequality on each cell, followed by integration, gives
\[
 \|(I_bf_j-f_j)'\|_2\le CbT^2/\sqrt a,
 \qquad
 \|I_bf_j-f_j\|_2\le Cb^2T^2/\sqrt a.
\]
No derivative jump lies inside a mesh cell. The sampling Sobolev
bound, equivalently \eqref{eq:fd-sampling-sobolev} or its displayed
Plancherel proof, yields
\[
 q(I_bf_j-f_j)\le Ca b^2T^4.
\]
Also \(|\int I_bf_j|\le Cb^2T^2\), so the exact mean correction
satisfies the same bound. Therefore
\[
 \|U_N-U\|_{\mathbb C^M\to\|\cdot\|_T}
   \le C\sqrt{Ma}\,bT^{5/2}
   \le CbT^3\sqrt{a\log(3T)}.
\]
Choose \(A\) also so that \(A\log2>3\). Together with
\eqref{eq:strong-inertia-fullnorm}, this makes the norm of the
pulled-back quadratic-form difference tend to zero:
\[
 \|U_N^*Q_TU_N-U^*Q_TU\|
 \le(\|U_N\|_T+\|U\|_T)\|U_N-U\|_T=o(1).
\]
The positive margin in \eqref{eq:strong-inertia-margin} persists.
Thus the range of \(U_N\) is an \(M\)-dimensional positive subspace
for \(Q_T\) in \(V_N\). Sylvester's law identifies its finite
matrix as \(T H_N-G_N\), and proves the stated inertia bound.
The dimension and prime range follow directly from
\(\dim V_N=2N2^N-2\) and \(N\le A\log T+1\).
\end{proof}

\subsection{Approximation of isolated simple coordinates}

\begin{theorem}[Localized spectral coordinates in logarithmic support]
\label{thm:fd-inertia-density}
For every \(0<\eta<c_{\rm MT}\) and \(\epsilon>0\), there are
constants \(A,\alpha,\beta,C,T_0>0\), with \(0<\alpha<\beta<1\),
such that for every \(T\ge T_0\) there is a set of simple ordinates
\(Z_T=\{\gamma_1,\ldots,\gamma_M\}\subset[\alpha T,\beta T]\)
with \(M\ge(c_{\rm MT}-\eta)n(T)\) and a map
\(U_N:\mathbb C^M\to V_N\), \(N=\lceil A\log T\rceil\), satisfying
\[
 \|U_N-E\|\le\epsilon,
 \qquad
 \|CU_N-U_N\Lambda\|\le\frac{C}{T^2\log T}.
\]
Here \(E\) sends the standard basis to the selected unit spectral
coordinates in \(\ell^2(\Gamma,m)\), and
\(\Lambda=\operatorname{diag}(1/\gamma_j)\); both norms are operator
norms into the Weil completion. The same maps give
\(\iota_N(T)\ge(c_{\rm MT}-\eta)n(T)\).
The dimension and prime range have the bounds of
Theorem~\ref{thm:strong-inertia-density}.
\end{theorem}

\begin{proof}
We give the approximation estimates explicitly. The proof of
Theorem~\ref{thm:two-thirds-survival}, before its heat continuation
step, supplies isolated simple ordinates with the following stronger
data. After restricting to finitely many adjacent dyadic intervals,
there are constants
\[
 0<\alpha<\beta<1,\qquad \delta>0,\qquad B>0
\]
and, for every sufficiently large \(T\), a set
\(Z_T=\{\gamma_1,\ldots,\gamma_M\}\subset[\alpha T,\beta T]\)
with
\begin{equation}\label{eq:fd-selected-data}
 \begin{split}
 M&\ge(c_{\rm MT}-\eta/2)n(T),\\
 |\gamma-\gamma_j|&\ge h:=\delta/\log T
     \quad(\gamma\ne\gamma_j),\\
 S_j:=\sum_{\gamma\ne\gamma_j}
           \frac{m_\gamma}{(\gamma-\gamma_j)^2}
     &\le B(\log T)^2.
 \end{split}
\end{equation}
The second line includes every other signed zero occurrence.
Here is the passage from the dyadic statement. Use the isolation
construction with a smaller error parameter on the finitely many
blocks ending at \(\beta T\) and beginning above \(\alpha T\).
Choose \(\beta\) sufficiently close to one and the lower endpoint
sufficiently small that the Riemann--von Mangoldt count loses less
than the allowed fraction. The proof already removes strips of width
\(\log T\) at each block endpoint, so adjacent selected blocks are
separated as well. There are only finitely many blocks, and their
heights are fixed positive multiples of \(T\). Their gap and
inverse-square constants may thus be replaced by common constants.
Changing from the heat coordinate \(x=2\gamma\) divides the gaps by
two and multiplies the inverse-square sums by four. This proves
\eqref{eq:fd-selected-data} with the displayed zero-ordinate
normalization.

Let \(\varphi\) be the convolution of eight copies of
\(4\mathbf1_{[-1/8,1/8]}\). It has integral one, support
\([-1,1]\), and square-integrable derivatives through order three.
Its Fourier transform in the convention used here is
\[
 K(v)=\left(\frac{\sin(v/8)}{v/8}\right)^8,
 \qquad K(0)=1.
\]
For every nonzero real \(v\),
\begin{equation}\label{eq:fd-bump-decay}
 |K(v)|\le\frac{64}{v^2},\qquad
 |vK(v)|\le\frac{512}{v^2}.
\end{equation}
Indeed, \(|K(v)|\le\min\{1,(8/|v|)^8\}\); split at
\(|v|=8\) to obtain both inequalities.

Put \(\ell=\log T\), choose \(N=\lceil A\ell\rceil\), and
let \(L=N/2\). For each selected ordinate define compact functions
\[
 g_j(x)=L^{-1}e^{i\gamma_jx}\varphi(x/L),\qquad
 f_j(x)=\frac1{i\gamma_j}\,g_j'(x).
\]
They are supported on \([-L,L]\), and \(f_j\) has integral zero.
Their Fourier transforms and compact primitives satisfy
\[
 \widehat f_j(\gamma)
    =\frac{\gamma}{\gamma_j}K(L(\gamma-\gamma_j)),
 \qquad C f_j=\frac{g_j}{\gamma_j}.
\]
The second identity follows either from the compact primitive
\(i\int_{-N}^x f_j\) or from division of the Fourier transform by
\(\gamma\).

Use the unitary spectral coordinates
\(f\mapsto(\sqrt{m_\gamma}\widehat f(\gamma))_\gamma\).
Let \(E:\mathbb C^M\to\mathcal H\) place the standard basis
at the selected simple coordinates, and let
\(U:\mathbb C^M\to\mathcal H\) have columns \(f_j\).
Define two matrices, omitting their selected diagonal entries,
\[
 \mathcal K_{\gamma j}=\sqrt{m_\gamma}
       K(L(\gamma-\gamma_j)),\qquad
 \mathcal L_{\gamma j}=\sqrt{m_\gamma}
       (\gamma-\gamma_j)K(L(\gamma-\gamma_j)).
\]
For every row ordinate, the selected ordinates on either side are
\(h\)-separated, and every retained distance is at least \(h\).
Consequently
\[
 \sum_{j:\gamma_j\ne\gamma}(\gamma-\gamma_j)^{-2}
       \le2h^{-2}\sum_{r\ge1}r^{-2}<4h^{-2}.
\]
Apply the weighted Schur test with row weights
\(\sqrt{m_\gamma}\) and column weights one. The row and column
bounds for \(\mathcal K\) are respectively
\(256/(L^2h^2)\) and \(64B\ell^2/L^2\). For
\(\mathcal L\) they are \(2048/(L^3h^2)\) and
\(512B\ell^2/L^3\). Thus
\begin{equation}\label{eq:fd-bump-matrix-bounds}
 \|\mathcal K\|\le\frac{128\sqrt B}{\delta}
                 \frac{\ell^2}{L^2},\qquad
 \|\mathcal L\|\le\frac{1024\sqrt B}{\delta}
                 \frac{\ell^2}{L^3}.
\end{equation}
The Schur test follows directly by applying Cauchy's inequality in
each row with these weights and then summing the column bounds.
In particular its use here includes every complementary multiplicity.

Let \(\Lambda=\operatorname{diag}(1/\gamma_j)\). The exact
Fourier formulas give
\[
 U-E=\mathcal K+\mathcal L\operatorname{diag}(1/\gamma_j),
 \qquad C U-U\Lambda
       =-\mathcal L\operatorname{diag}(1/\gamma_j^2).
\]
Since \(L\ge A\ell/2\), these imply
\begin{align}
 \|U-E\|&\le\frac{512\sqrt B}{\delta A^2}
       +\frac{8192\sqrt B}{\delta\alpha A^3T\ell},
       \label{eq:fd-bump-isometry}\\
 \|CU-U\Lambda\|&\le
       \frac{8192\sqrt B}{\delta\alpha^2A^3T^2\ell}.
       \label{eq:fd-bump-residual}
\end{align}
These estimates do not use a lower bound for a zeta derivative.

Put \(g=1/\beta-1>0\), \(M_0=1/\alpha-1\), and
\(\varepsilon_0=\min\{1/4,g/(12M_0)\}\).
Choose \(A\ge8\) sufficiently large and then \(T\) sufficiently
large that \(\|U-E\|\le\min(\varepsilon_0,\epsilon/2)\).
It follows that
\(\|U^*U-I\|\le3\varepsilon_0\).
For a unit vector \(a\in\mathbb C^M\), write
\(R=CU-U\Lambda\). Taking real parts in the exact identity
for the quadratic form gives
\begin{align*}
 \langle(C-T^{-1})Ua,Ua\rangle
 &\ge\frac gT-
       \|U^*U-I\|\frac{M_0}{T}-\|U\|\|R\|\\
 &\ge\frac{g}{2T}
\end{align*}
for all sufficiently large \(T\), by
\eqref{eq:fd-bump-residual}. Also \(U\) is injective, since it
is within \(1/4\) of the isometry \(E\).

It remains to place this whole positive subspace in the prescribed
finite piecewise-linear space. We include the needed norm estimates.
For every \(f\in H^1(\mathbb R)\) supported on \([-N,N]\),
\(N\ge1\), we have
\begin{equation}\label{eq:fd-sampling-sobolev}
 \|f\|_q\le50N\|f\|_{H^1}.
\end{equation}
Here are explicit constants. The completion estimates in
Lemma~\ref{lem:fd-multiplicity-log}, now at an arbitrary real ordinate
\(v\), give
\[
 \Re\frac{\xi'}{\xi}(2+iv)
 <\frac{13}{5}+\frac12\log(|v|+2)
 <4\log(|v|+3).
\]
Each zero in \([k,k+1)\) contributes at least \(3/5\) to the
positive logarithmic-derivative sum at \(v=k+1/2\). The counted
zeros in this interval therefore number at most
\(10\log(|k|+4)\). On a unit interval,
\(\sup|F|^2\le2\int(|F|^2+|F'|^2)\): choose a point where
\(|F|^2\) is no larger than its integral and integrate \(F'\)
from that point. For \(F=\widehat f\), summing gives
\[
 q(f,f)\le60\int_{\mathbb R}(1+v^2)
                   (|F(v)|^2+|F'(v)|^2)\,dv,
\]
where we used \(\log(|v|+5)\le3(1+v^2)\).
Plancherel identifies the integral as
\[
 2\pi\{\|f\|_2^2+\|f'\|_2^2+
                  \|xf\|_2^2+\|(xf)'\|_2^2\}
 \le6\pi(1+N^2)\|f\|_{H^1}^2.
\]
Thus \(q(f,f)\le720\pi N^2\|f\|_{H^1}^2
<2500N^2\|f\|_{H^1}^2\), proving
\eqref{eq:fd-sampling-sobolev}.

Write \(b=2^{-N}\), and let \(I_b f\) be linear interpolation
on the mesh in \([-N,N]\). The cellwise Poincar\'e inequality gives
\[
 \|(I_bf-f)'\|_2\le b\|f''\|_2,\qquad
 \|I_bf-f\|_2\le b^2\|f''\|_2.
\]
Let \(v_0(x)=(1-|x|)_+\), which belongs to the mesh space and
has integral one. Correct the mean by defining
\[
 J_bf=I_bf-v_0\int_{-N}^{N}I_bf.
\]
For mean-zero \(f\),
\[
 \left|\int I_bf\right|\le\sqrt{2N}\,b^2\|f''\|_2.
\]
Since \(b\sqrt N\le1\) and
\(\|v_0\|_{H^1}=\sqrt{8/3}\), the preceding three estimates imply
\begin{equation}\label{eq:fd-interpolation-h1}
 \|J_bf-f\|_{H^1}\le5b\|f''\|_2.
\end{equation}
The corrected interpolant belongs to \(V_N\).

For our functions, direct differentiation gives
\[
 f_j''(x)=e^{i\gamma_jx}
 \left[-\frac{\gamma_j^2}{L}\varphi(x/L)
       +\frac{3i\gamma_j}{L^2}\varphi'(x/L)
       +\frac3{L^3}\varphi''(x/L)
       +\frac1{i\gamma_jL^4}\varphi'''(x/L)\right].
\]
For large \(T\), \(1\le\gamma_j\le T\) and \(L\ge1\),
so
\[
 \|f_j''\|_2\le C_\varphi T^2/\sqrt L,\qquad
 C_\varphi=\|\varphi\|_2+3\|\varphi'\|_2+
             3\|\varphi''\|_2+\|\varphi'''\|_2\le1458.
\]
For the last bound, each original box has \(L^1\) norm one and
\(L^2\) norm two, and its distributional derivative has total
variation eight. Young's convolution inequality therefore gives
\(\|\varphi^{(r)}\|_2\le2\cdot8^r\) for \(0\le r\le3\).
Let \(U_N\) have columns \(J_bf_j\). Combining
\eqref{eq:fd-sampling-sobolev} and
\eqref{eq:fd-interpolation-h1}, then using
\(M\le C T\log T\), gives
\begin{align*}
 \|U_N-U\|
 &\le\sqrt M\max_j\|J_bf_j-f_j\|_q\\
 &\le C_\eta\,2^{-N}T^{5/2}\log T
 \le C_\eta T^{5/2-A\log2}\log T=o(T^{-1}).
\end{align*}
Here \(A\ge8\) makes the final exponent strictly less than
\(-1\). The form difference is bounded by
\[
 \|U_N^*(C-T^{-1})U_N-U^*(C-T^{-1})U\|
 \le(\|C\|+T^{-1})
       (\|U_N\|+\|U\|)\|U_N-U\|=o(T^{-1}).
\]
Consequently the form \(q(Cf,f)-T^{-1}q(f,f)\) is positive
definite on the \(M\)-dimensional range of \(U_N\) for all large
\(T\). The finite variational principle gives
\(\iota_N(T)\ge M\). The same interpolation estimate is
\(o(T^{-2}/\log T)\), since \(A\ge8\). Thus
\(\|U_N-E\|\le\epsilon\) for large \(T\), and
\[
 \|CU_N-U_N\Lambda\|
 \le\|CU-U\Lambda\|+(\|C\|+\|\Lambda\|)\|U_N-U\|
 \le C/(T^2\log T).
\]
This proves the coordinate and residual assertions. Finally, \(N\le A\log T+1\) and the explicit arithmetic formula
prove the stated prime range and matrix dimension.
\end{proof}

This estimate applies at every sufficiently large height and fixes
the support and mesh needed to realize the proved density in the
arithmetic matrices. The integer deficit
\(n(T)-\iota_N(T)\) is still permitted to be of order \(n(T)\).
Its vanishing at unbounded heights is the full simplicity gate in
Proposition~\ref{prop:fd-inertia-complete}; the approximation estimates
above do not establish that vanishing for the complementary set.

\subsection{The exact trace integral and the absolute integer scale}

Define the nonnegative integer deficit
\[
 A_N(T)=n(T)-\iota_N(T).
\]
The trace and inertia certificates are related by the exact identity
\begin{equation}\label{eq:fd-trace-inertia-integral}
 D_{N,k}=4k\int_0^\infty A_N(t)t^{-2k-1}\,dt,
 \qquad k\ge1.
\end{equation}
To prove it, parity pairs the nonzero eigenvalues of \(C_N\)
as \(\pm\lambda_{N,j}\), with \(\lambda_{N,j}>0\). For a
positive number \(\lambda\),
\[
 \lambda^{2k}=2k\int_{1/\lambda}^{\infty}t^{-2k-1}\,dt.
\]
Apply this identity first to the counted reciprocal ordinates in
\(S_k\), and then to the positive compression eigenvalues in
\(\operatorname{tr}(C_N^{2k})\). Both integrals converge:
the counting functions vanish near zero and grow at most as
\(O(t\log(t+3))\) at infinity. Subtraction proves
\eqref{eq:fd-trace-inertia-integral}; choices at the isolated
threshold points do not change the integrals. At zero-free heights,
\(A_N(T)\) decreases to the multiplicity excess
\(n(T)-n_{\mathrm{dist}}(T)\).

There is also an exact finite-dimensional comparison with the
positive-kernel replacement of
Section~\ref{sec:fd-positive-surgery}. Its prescribed simple
ordinates satisfy \(0<u<v\): the replacement removes the pair
\(\pm u\) and adds one occurrence at each of \(\pm v\).
Use the same compact mean-zero basis to form both Gram matrices.
For \(f,g\) in its span, write \(F=\widehat f\), \(G=\widehat g\)
and \(\mathcal P f=i\int_{-N}^{x}f\). The exact change of the
form \(T H_N-G_N\) is
\begin{align}
 &\widetilde q(T\mathcal P f-f,g)-q(T\mathcal P f-f,g)
 \notag\\
 &\quad=\sum_{\sigma=\pm1}\left[
  \left(\frac{T}{\sigma v}-1\right)
            F(\sigma v)\overline{G(\sigma v)}
 -\left(\frac{T}{\sigma u}-1\right)
            F(\sigma u)\overline{G(\sigma u)}\right].
 \label{eq:fd-inertia-surgery}
\end{align}
This follows by applying the already proved four-evaluation change
of \(q\) and using
\(\widehat{\mathcal P f}(\gamma)=F(\gamma)/\gamma\).
It compares forms on their common test functions and makes no
identification of the completed Hilbert spaces.

For \(T>v\), the terms at \(+v\) and \(-u\) have positive
coefficients, while those at \(-v\) and \(+u\) have negative
coefficients. Thus the finite matrix change is a sum of at most two
positive rank-one forms minus at most two positive rank-one forms.
The variational principle, applied to the kernels of these forms,
gives
\[
 |\widetilde\iota_N(T)-\iota_N(T)|\le2.
\]
For zero-free \(T>v\), the total counted zero numbers are equal,
whereas the distinct count for the replacement is smaller by one.
The eventual inertia counts therefore give
\[
 \lim_N\widetilde A_N(T)=\lim_N A_N(T)+1.
\]
The sampling construction and compact projections also apply to the
replacement. Indeed, its unit-interval bounds are preserved by the
finite change. A compactly supported nonzero function has a Fourier
transform of exponential type and therefore only \(O(R)\) zeros in
a disk of radius \(R\), by Jensen's formula centered at a nonzero
value. The unchanged positive-density set of simple ordinates has
order \(R\log R\), so the modified form is strictly positive on
such test functions. To verify density, an orthogonal vector would
define a tempered distribution supported on the modified ordinates
whose inverse Fourier transform annihilates all compact mean-zero
tests. That inverse transform must be constant, so the original
distribution is supported at zero. There is no zero ordinate, hence
the orthogonal vector vanishes. For support in \([-N,N]\), the
four-evaluation update and Cauchy's inequality give
\[
 \widetilde q(f)\le q(f)+|F(v)|^2+|F(-v)|^2
 \le q(f)+4N\|f\|_2^2.
\]
The sharper bound in the proof of
\eqref{eq:fd-sampling-sobolev} was
\(q(f)\le720\pi N^2\|f\|_{H^1}^2\). Since
\(720\pi+4<2500\) and \(N\ge1\), the same displayed
\(50N\) norm bound holds for \(\widetilde q\).
Mean-corrected interpolation therefore shows that the union of the
finite spaces is dense as well. This justifies the eventual inertia count
for the replacement. In particular the positive-density lower estimate
is preserved after changing its error parameter, since the finite
change is bounded by two and \(n(T)\to\infty\).

Equations \eqref{eq:fd-trace-inertia-integral} and
\eqref{eq:fd-inertia-surgery} specify the absolute scale of the
remaining estimate. The positive-density approximation controls a
fraction of the integer deficit. The displayed finite replacement
preserves that fractional conclusion while changing its eventual
integer value by one, and its change on the arithmetic side is exactly
the oscillatory Weil correction already computed. An estimate proving
zero deficit must therefore control this integer scale for the actual
prime form.
